\section{A normalized conditional Wilson tail in a fixed box}
\label{fbw:section}

This section controls the entire compact Wilson integral for a fixed
box with small boundary data. The constants depend on the box. It
does not identify this box gauge with the nonlinear constrained
coordinates of the homogeneous renormalization step.

Fix a cubical lattice box \(\Omega=\{0,\ldots,R\}^d\), \(d=4\),
\(R\geq2\), with its full cubical cell structure. Let \(E_\partial\)
be the edges in its topological boundary and \(E_\circ\) the remaining
edges. The boundary edge graph is connected. Fix
\(G=SU(N)\), with faithful normalized trace \(\tau=\operatorname{Tr}/N\),
and put
\[
 h(U)=1-\operatorname{Re}\tau(U)
     =\tfrac12\norm{U-I}_\tau^2,\qquad
 S_V(U)=\sum_{p\subset\Omega}h(U_p).
\]
Boundary edge variables \(V\in G^{E_\partial}\) are prescribed.
All remaining edges have normalized Haar measure. Gauge transformations
are fixed to the identity at every boundary vertex.

Choose a spanning forest \(\mathcal F\) of the full vertex set with
exactly one boundary vertex in every component, and all interior
vertices included. Its edges lie in \(E_\circ\). Such a forest exists:
start with a spanning tree of the connected boundary graph, extend
it to a spanning tree of the whole graph, and remove its boundary
tree edges. There is a unique allowed gauge transformation making
every forest edge the identity. Write \(y\in\mathcal Y=G^{E_\circ
\setminus\mathcal F}\) for the remaining edge variables in that gauge.
The number of real coordinates is
\[
 n=(\dim G)\,|E_\circ\setminus\mathcal F|.
\]
For this box \(n>0\). Fix an orthonormal real Lie algebra basis and
write \(y=\exp(ix)\) componentwise on a neighborhood of the identity.
Here \(\norm{x}_\infty=\max_e\norm{x_e}_\tau\) uses the
adjoint-invariant Lie algebra norm on each edge.

\begin{theorem}[Full conditional tail at fixed geometry]
\label{fbw:theorem}
There are a neighborhood \(\mathcal V\) of the identity boundary
configuration, constants \(g_0,r,c,C>0\), and a real analytic
function \(x_*(V)\) on \(\mathcal V\), with \(x_*(I)=0\), such that
the following holds. The forest-gauge action has the unique global
minimum \(y_*(V)=\exp(ix_*(V))\). For
\(0<g\leq g_0\) and \(1\leq T\leq r/g\), its normalized Wilson
measure satisfies
\begin{equation}
 \frac{\displaystyle
    \int_{\mathcal Y\setminus W_{V,g,T}}
           e^{-S_V(y)/g^2}\,dy}
      {\displaystyle\int_{\mathcal Y}e^{-S_V(y)/g^2}\,dy}
 \leq C e^{-cT^2},
 \qquad
 W_{V,g,T}=\{\exp(ix):\norm{x-x_*(V)}_\infty\leq gT\}.
 \label{fbw:tail}
\end{equation}
The displayed window is contained in a single fixed exponential
coordinate neighborhood. The same ratio is the probability of the
corresponding gauge-fixed event under the original, ungauged box
Wilson measure.
\end{theorem}

\begin{proof}
We give the global normalization argument, including the gauge and
topological points.

\emph{Exact forest gauge.}
Orient each forest component away from its boundary root. Recursively
solving the transformation equation along its edges gives the unique
interior gauge values. Changing from the original forest edge variables
to these gauge values is triangular; each step is a left or right
translation, or inversion, on \(G\), hence preserves normalized Haar
measure. For fixed gauge values the other edge transformations preserve
Haar independently. The change of variables is therefore an exact
product-Haar identity. A gauge-invariant action and an event defined
through the unique forest representative are independent of the
interior gauge variables. Their integrals over those variables equal
one. No determinant or volume factor is discarded. Boundary edges
are unchanged even when \(V\ne I\); this is why a forest rooted at
all boundary vertices is used.

\emph{The unique zero and its Hessian.}
At \(V=I\), zero action means every plaquette holonomy is \(I\).
The cubical box is simply connected, so a flat edge field is a pure
gauge: parallel transport from a fixed vertex is independent of the
path, since elementary plaquette moves generate all path homotopies.
Identity boundary edges force the vertex gauge to be constant on the
connected boundary; normalize that constant to \(I\). The field
thus has an allowed gauge transformation to the identity. Uniqueness
of the forest representative gives \(y=I\). Consequently the
continuous nonnegative function \(S_I\) has exactly one zero on the
compact space \(\mathcal Y\).

For a Lie algebra edge field \(a\), zero on boundary and forest edges,
the quadratic part is
\[
 \tfrac12\sum_{p\subset\Omega}\norm{(da)_p}_\tau^2.
\]
If it vanishes, the elementary cellular Poincare lemma on a box gives
\(a=d\phi\). Boundary connectedness lets us choose \(\phi=0\) at
every boundary vertex. Vanishing on each forest edge then gives
\(\phi=0\) at every vertex. Hence \(a=0\). The Hessian in the \(n\)
remaining coordinates is positive definite. This is a finite
dimensional assertion; no uniform lower eigenvalue as \(R\) varies
is asserted.

\emph{Uniform local well and global gap.}
The analytic implicit-function theorem supplies a critical point
\(x_*(V)\) near zero for \(V\) close to \(I\). Choose a coordinate
ball about zero on which the Hessian at all such \(V\) is uniformly
positive. On a smaller fixed ball about \(x_*(V)\), Taylor's theorem
gives constants \(c_0,C_0>0\) with
\begin{equation}
 c_0|x-x_*(V)|^2\leq S_V(\exp(ix))-S_V(y_*(V))
       \leq C_0|x-x_*(V)|^2.
 \label{fbw:quadratic}
\end{equation}
The balls can be chosen uniformly after shrinking \(\mathcal V\).
To see a global gap, first choose a still smaller ball about zero.
On the compact complement its \(V=I\) action has a strictly positive
minimum. Uniform continuity of \(S_V\), and \(S_V(y_*(V))\to0\),
preserve half this gap for sufficiently small boundary data.
Together with strict convexity on the coordinate ball this proves
uniqueness of the global minimum and yields a fixed chart region
\(\mathcal U\subset\mathcal Y\) and \(\eta>0\) for which
\begin{equation}
 S_V(y)\geq S_V(y_*(V))+\eta\quad(y\notin\mathcal U).
 \label{fbw:gap}
\end{equation}
Choose \(\mathcal U\) small enough that the lower quadratic bound in
\eqref{fbw:quadratic} holds on all of it, and large enough that a
fixed ball about every \(x_*(V)\) lies in it. These choices are
compatible by first making the convexity ball larger than the
compact-complement ball and then shrinking \(\mathcal V\).
The product Haar coordinate density on its closure is bounded above
and below by positive constants.

\emph{Denominator and local numerator.}
Writing \(S_{\min}=S_V(y_*(V))\), integrate over the coordinate
ball \(|x-x_*|\leq g\), reducing \(g_0\) as needed:
\begin{equation}
 Z_V(g)\geq c_1 g^n e^{-S_{\min}/g^2}.
 \label{fbw:denominator}
\end{equation}
Within \(\mathcal U\), the complement of the window forces
\(|x-x_*|\geq gT\). The lower quadratic bound and the upper Haar
density bound give, for \(T\geq1\),
\begin{align}
 \int_{\mathcal U\setminus W_{V,g,T}}e^{-S_V/g^2}\,dy
 &\leq C_1 e^{-S_{\min}/g^2}
       \int_{|z|\geq gT}e^{-c_0|z|^2/g^2}\,dz \notag\\
 &\leq C_2 g^n e^{-S_{\min}/g^2}e^{-c_2T^2}.
 \label{fbw:localnumerator}
\end{align}
The elementary Gaussian tail estimate absorbs its finite dimensional
polynomial factor by reducing \(c_2>0\).

\emph{The rest of the compact group integral.}
The remaining numerator is at most
\(e^{-(S_{\min}+\eta)/g^2}\), since product Haar mass is one.
After \eqref{fbw:denominator} its contribution is at most
\(C_3g^{-n}e^{-\eta/g^2}\). Since
\[
 \sup_{0<g\leq g_0}g^{-n}e^{-\eta/(2g^2)}<\infty,
\]
and \(T\leq r/g\), choosing \(r\) small and
\(c\leq \min\{c_2,\eta/(2r^2)\}\) bounds this contribution by
\(C_4e^{-cT^2}\). Reduce \(r\) also to keep the entire sup-norm
window inside the coordinate chart. Combining the two numerators
with the denominator proves \eqref{fbw:tail}.
\end{proof}

\begin{corollary}[Logarithmic windows]
\label{fbw:logwindow}
For fixed \(a>1/2\), take \(T(g)=(\log(1/g))^a\). For every
\(m>0\) there are \(g_m,C_m>0\) such that
\[
 \mu_{V,g}(\mathcal Y\setminus W_{V,g,T(g)})\leq C_m g^m
 \quad(0<g\leq g_m,\ V\in\mathcal V).
\]
\end{corollary}
\begin{proof}
Here \(gT(g)\to0\) and \(T(g)^2/\log(1/g)\to\infty\).
Apply \eqref{fbw:tail}.
\end{proof}

\paragraph{Conditional interpretation.}
Embed the box in a larger lattice without local periodic
identifications. A plaquette containing an edge of \(E_\circ\)
lies entirely inside the box: an edge not in the boundary has all
its transverse coordinates strictly inside the corresponding
intervals. Thus, after the exterior and boundary edges are fixed,
the part of the bare Wilson action depending on the integrated edges
is exactly \(S_V\), up to an additive constant. Such constants cancel
in the normalized law. The theorem is therefore a statement about
the actual conditional bare Wilson measure whenever its boundary
field belongs to \(\mathcal V\). A prescribed boundary gauge
transformation bringing the data into \(\mathcal V\) also permits
the theorem after transforming the event consistently.

\paragraph{The precise interface left open.}
This result controls a fixed-box forest-coordinate window about
the unconstrained conditional minimizer. It does not control an
arbitrary boundary field, a box growing with the cutoff, a
nonlinear coarse-average fiber, or a mixed-history effective action
whose global stability has not been proved. In particular an
unconditional probability bound cannot simply be conditioned on a
coarse-map value with no lower bound for the corresponding fiber
normalization. The powers \(g^n\) cancel here because both integrals
are estimated in the same compact conditional model. Extending
that cancellation to the actual constrained local map is a
remaining analytic task.
